\chapter{Estimation}\label{ch:qm:estimation}

Un chercheur présente une stratégie~: sur cinq ans, son P\&L quotidien moyen atteint 4.0~points de base, avec une
\omterm{def:qm:estimation:estimator}{erreur type} de 1.2, une statistique $t$ de 3.4, nettement significative. Un collègue remarque que
chaque position est conservée pendant cinq jours, de sorte que deux jours consécutifs partagent quatre cinquièmes de leurs
positions, et recalcule l'\omterm{def:qm:estimation:estimator}{erreur type} en tenant compte de l'autocorrélation que cela crée~: 2.3.
La statistique $t$ tombe à 1.7, et la stratégie ne se distingue plus du hasard. Rien
n'était faux dans la moyenne~; la formule de son incertitude supposait des jours indépendants. Ce chapitre
porte sur les \omterm{def:qm:estimation:estimator}{estimateurs} et leur incertitude~: le maximum de vraisemblance et ses propriétés asymptotiques, ce qui se passe
lorsque le modèle est faux, les \omterm{def:qm:estimation:m}{M-estimateurs} et les moments, la \omterm{met:qm:estimation:delta}{méthode delta} qui donne au \omterm{def:qm:estimation:sharpe}{ratio de Sharpe} son
\omterm{def:qm:estimation:estimator}{erreur type}, et les variances convergentes en présence d'hétéroscédasticité et d'autocorrélation (\omterm{def:qm:estimation:hac}{estimateurs HAC}) qui ont corrigé
la statistique $t$ du chercheur.

\section{Maximum de vraisemblance}

\begin{definition}[Estimateur, sans biais, convergent, erreur type]\label{def:qm:estimation:estimator}
Un \emph{estimateur}\index{estimateur} $\hat\theta_n$ d'un paramètre $\theta_0$ est une fonction des données $X_1, \dots, X_n$. Il est un
\emph{estimateur sans biais}\index{estimateur sans biais} si $\E[\hat\theta_n] = \theta_0$, et un \emph{estimateur
convergent}\index{estimateur convergent} si $\hat\theta_n \xrightarrow{\P} \theta_0$ lorsque $n \to \infty$. Son \emph{erreur
type}\index{erreur type} $\mathrm{se}(\hat\theta_n)$ est (une estimation de) son écart type.
\end{definition}

La variance empirique de diviseur $n$ est biaisée et convergente~; avec le diviseur $n - 1$, elle est sans biais. La convergence
est ce qui compte avec les tailles d'échantillon de la finance, et il en va de même de l'\omterm{def:qm:estimation:estimator}{erreur type}~: une estimation sans \omterm{def:qm:estimation:estimator}{erreur type}
n'est qu'une anecdote.

\begin{definition}[Fonction de vraisemblance, estimateur du maximum de vraisemblance, score, information de Fisher]\label{def:qm:estimation:mle}
Pour des observations indépendantes de densité $f(x; \theta)$, la \emph{fonction de vraisemblance}\index{fonction de vraisemblance} est
$\theta \mapsto \prod_if(X_i; \theta)$ et la log-vraisemblance $\ell_n(\theta) = \sum_i\ln f(X_i; \theta)$. L'\emph{estimateur du maximum
de vraisemblance}\index{estimateur du maximum de vraisemblance} (MLE) la maximise. La \emph{fonction score}\index{fonction score}
est $\nabla_\theta\ln f(X; \theta)$, et l'\emph{information de Fisher}\index{information de Fisher} est $\mathcal I(\theta) =
\E[\nabla\ln f\,\nabla\ln f^\top] = -\E[\nabla^2\ln f]$ par observation.
\end{definition}

\begin{theorem}[Asymptotique du maximum de vraisemblance]\label{thm:qm:estimation:mle}
Sous des conditions de régularité (identifiabilité, différentiabilité, vraie valeur à l'intérieur de l'espace des paramètres), le MLE est convergent et vérifie~:
\[
\sqrt n(\hat\theta_n - \theta_0) \xrightarrow{d} \mathcal N\bigl(0, \mathcal I(\theta_0)^{-1}\bigr).
\]
Aucun \omterm{def:qm:estimation:estimator}{estimateur sans biais} n'a une variance inférieure à $\mathcal I(\theta_0)^{-1}/n$ (la borne de Cramér--Rao), de sorte que le MLE est
asymptotiquement efficace.
\end{theorem}

\begin{proof}[Démonstration partielle]
Le score en $\theta_0$ est d'espérance nulle et de variance $\mathcal I$, donc $n^{-1/2}\nabla\ell_n(\theta_0) \to \mathcal N(0, \mathcal I)$ par le
théorème central limite. En développant $0 = \nabla\ell_n(\hat\theta_n) = \nabla\ell_n(\theta_0) + \nabla^2\ell_n(\tilde\theta)(\hat\theta_n - \theta_0)$ et en utilisant
$-n^{-1}\nabla^2\ell_n \to \mathcal I$, on obtient $\sqrt n(\hat\theta_n - \theta_0) = \mathcal I^{-1}n^{-1/2}\nabla\ell_n(\theta_0) + o_{\P}(1)$. La convergence
et la borne de Cramér--Rao sont traitées dans Casella et Berger (2002).
\end{proof}

\begin{example}[Le taux d'arrivées]\label{ex:qm:estimation:rate}
Pour $n$ temps d'attente exponentiels de taux $\lambda$, $\ell_n(\lambda) = n\ln\lambda - \lambda\sum t_i$, donc $\hat\lambda = n/\sum t_i$, $\mathcal I(\lambda) =
1/\lambda^2$, et $\mathrm{se}(\hat\lambda) = \hat\lambda/\sqrt n$. Cent intervalles totalisant 50 secondes donnent $2 \pm 0.2$ arrivées par seconde.
Les ajustements de Hawkes du chapitre~7 relèvent du même calcul, avec une vraisemblance plus complexe.
\end{example}

\section{Erreur de spécification et sandwich}

Les modèles sont des approximations~; la question est de savoir ce qu'estime le maximum de vraisemblance lorsque la densité est fausse.

\begin{definition}[Divergence de Kullback--Leibler, quasi-maximum de vraisemblance]\label{def:qm:estimation:kl}
La \emph{divergence de Kullback--Leibler}\index{divergence de Kullback--Leibler} d'une densité $f$ par rapport à la densité vraie $p$
est $D_{\mathrm{KL}}(p\,\|\,f) = \E_p[\ln(p(X)/f(X))] \ge 0$, nulle seulement si $f = p$. Le \emph{quasi-maximum
de vraisemblance}\index{quasi-maximum de vraisemblance} maximise une vraisemblance que l'on ne croit pas être la vraie.
\end{definition}

\begin{proposition}[Ce qu'estime une vraisemblance fausse]\label{prop:qm:estimation:qmle}
Sous des conditions de régularité, le quasi-MLE converge vers la valeur pseudo-vraie $\theta^*$ qui minimise $D_{\mathrm{KL}}(p\,\|\,f_\theta)$,
et $\sqrt n(\hat\theta_n - \theta^*) \xrightarrow{d} \mathcal N(0, A^{-1}BA^{-1})$ avec $A = -\E[\nabla^2\ln f_{\theta^*}]$ et $B = \Var(\nabla\ln
f_{\theta^*})$.
\end{proposition}

\begin{proof}
$n^{-1}\ell_n(\theta) \to \E_p[\ln f_\theta] = \E_p[\ln p] - D_{\mathrm{KL}}(p\,\|\,f_\theta)$, maximisé en $\theta^*$. Le développement de la
démonstration précédente tient, en remplaçant $\mathcal I$ par $A$ dans la hessienne et par $B$ dans la variance du score~; elles ne coïncident que lorsque
le modèle est juste (l'égalité de l'information).
\end{proof}

\begin{definition}[Variance sandwich]\label{def:qm:estimation:sandwich}
La \emph{variance sandwich}\index{variance sandwich} d'un \omterm{def:qm:estimation:estimator}{estimateur} défini par une équation d'estimation est $A^{-1}BA^{-1}$~:
le pain $A$, la dérivée espérée de l'équation d'estimation, autour de la garniture $B$, la variance de l'équation
elle-même. Estimée par la hessienne empirique et le produit extérieur des scores par observation, c'est la variance robuste
(de Huber--White).
\end{definition}

Ajustons un modèle normal à 2\,000 tirages d'une loi de Student $t$ à 6 degrés de liberté. L'estimation du logarithme de l'écart
type reste convergente vers le logarithme de l'écart type, mais l'\omterm{def:qm:estimation:estimator}{erreur type} par la hessienne, 0.0158, suppose des
queues normales~; le sandwich donne 0.0230, et sur 4\,000 échantillons simulés la dispersion réelle est de 0.0237
(\cref{fig:qm:estimation:sandwich}). Avec des queues encore plus épaisses (4 degrés de liberté, moment d'ordre quatre infini), la
variance de la variance empirique n'existe pas, et aucune \omterm{def:qm:estimation:estimator}{erreur type} n'est juste.

\begin{omfigure}
\begin{tikzpicture}
\begin{axis}[width=11cm,height=4.8cm,
  xlabel={$\ln\hat\sigma - $ sa moyenne},ylabel={densité},
  tick label style={font=\small},label style={font=\small},
  xmin=-0.1,xmax=0.1,ymin=0,ymax=27,grid=major,grid style={gray!25},
  xtick={-0.1,-0.05,0,0.05,0.1},xticklabel style={/pgf/number format/fixed},scaled x ticks=false,
  legend columns=3,legend style={font=\footnotesize,at={(0.5,-0.3)},anchor=north,draw=none,fill=none,/tikz/every even column/.append style={column sep=8pt}},
  table/col sep=comma]
\addplot[black!60,thick,const plot mark mid] table[x=x,y=hist]{figdata/methods/11-estimation/sandwich.csv};
\addplot[omThm,very thick,dashed] table[x=x,y=hessian]{figdata/methods/11-estimation/sandwich.csv};
\addplot[omDef,very thick] table[x=x,y=sandwich]{figdata/methods/11-estimation/sandwich.csv};
\legend{4\,000 échantillons,hessienne (modèle normal),sandwich}
\end{axis}
\end{tikzpicture}

\omcaption{Le logarithme de l'écart type estimé par une vraisemblance normale à partir de 2\,000 observations de Student-$t_6$~: sa
distribution d'échantillonnage réelle (en escalier) face aux approximations normales avec l'\omterm{def:qm:estimation:estimator}{erreur type} par la hessienne (tirets), trop étroite,
et avec le sandwich (trait plein). Données~: le tutoriel du chapitre, avec graine fixée.}\label{fig:qm:estimation:sandwich}
\end{omfigure}

\section{Estimateurs M et méthode des moments}

\begin{definition}[M-estimateur, méthode des moments, méthode des moments généralisée]\label{def:qm:estimation:m}
Un \emph{M-estimateur}\index{M-estimateur} maximise (ou annule la dérivée de) une moyenne empirique $n^{-1}\sum_im(X_i; \theta)$~;
le maximum de vraisemblance prend $m = \ln f$, les moindres carrés $m = -(y - x^\top\theta)^2$. La \emph{méthode des
moments}\index{méthode des moments} résout $n^{-1}\sum_ig(X_i) = \E_\theta[g(X)]$ pour autant de moments que de paramètres. La
\emph{méthode des moments généralisée}\index{méthode des moments généralisée} (GMM) minimise $\bar g_n(\theta)^\top W\bar g_n(\theta)$ pour
$\bar g_n(\theta) = n^{-1}\sum_ig(X_i; \theta)$, avec plus de conditions de moments que de paramètres et une matrice de pondération $W$.
\end{definition}

Les trois admettent les propriétés asymptotiques en sandwich de la \cref{prop:qm:estimation:qmle}, avec $A$ la dérivée des équations
d'estimation et $B$ leur variance de long terme~; la GMM est efficace lorsque $W$ est l'inverse de cette variance (Hansen, 1982),
et l'objectif minimisé teste alors les conditions de suridentification. Pour un \omterm{def:qm:stochastic-differential-equations:ou}{processus d'Ornstein--Uhlenbeck} échantillonné
quotidiennement, la \omterm{def:qm:estimation:m}{méthode des moments} appliquée à la première autocorrélation, $\hat\kappa = -\ln\hat\rho(1)/\Delta t$, coïncide avec le
maximum de vraisemblance conditionnel de l'autorégression (chapitre~17), et les deux héritent de son biais en petit échantillon.

\section{La méthode delta et le ratio de Sharpe}

\begin{method}[Méthode delta]\label{met:qm:estimation:delta}
Si $\sqrt n(\hat\theta_n - \theta_0) \xrightarrow{d} \mathcal N(0, V)$ et si $g$ est différentiable en $\theta_0$, alors $\sqrt n(g(\hat\theta_n) - g(\theta_0))
\xrightarrow{d} \mathcal N(0, \nabla g^\top V\nabla g)$. La \emph{méthode delta}\index{méthode delta} linéarise une fonction régulière d'un
\omterm{def:qm:estimation:estimator}{estimateur} et propage sa variance.
\end{method}

\begin{definition}[Ratio de Sharpe]\label{def:qm:estimation:sharpe}
Le \emph{ratio de Sharpe}\index{ratio de Sharpe} d'une stratégie est son rendement en excès espéré par unité d'écart type,
$\mathrm{SR} = \mu/\sigma$, par période~; il est annualisé par $\sqrt{\text{nombre de périodes par an}}$ lorsque les rendements sont indépendants d'une période à l'autre. Son
\omterm{def:qm:estimation:estimator}{estimateur} est $\widehat{\mathrm{SR}} = \bar x/s$.
\end{definition}

\begin{proposition}[Erreur type du ratio de Sharpe]\label{prop:qm:estimation:srse}
Pour des rendements normaux indépendants, $\sqrt n(\widehat{\mathrm{SR}} - \mathrm{SR}) \xrightarrow{d} \mathcal N(0, 1 + \tfrac12\mathrm{SR}^2)$ par période.
\end{proposition}

\begin{proof}
$(\bar x, s^2)$ sont indépendants, de variances asymptotiques $\sigma^2$ et $2\sigma^4$. Le gradient de $\mu/\sqrt v$ vaut $(1/\sigma,
-\mu/(2\sigma^3))$, si bien que la \omterm{met:qm:estimation:delta}{méthode delta} donne $1 + \mu^2\cdot2\sigma^4/(4\sigma^6) = 1 + \tfrac12\mathrm{SR}^2$.
\end{proof}

Annualisée, l'\omterm{def:qm:estimation:estimator}{erreur type} vaut environ $\sqrt{(1 + \mathrm{SR}_1^2/2)/\text{années}}$, puisque le \omterm{def:qm:estimation:sharpe}{ratio de Sharpe} journalier $\mathrm{SR}_1$ est petit~:
cinq années permettent d'estimer un \omterm{def:qm:estimation:sharpe}{ratio de Sharpe} annuel à $\pm 0.45$, quelle que soit sa valeur. La stratégie du chercheur a $\widehat{\mathrm{SR}} = 1.51$
avec cette \omterm{def:qm:estimation:estimator}{erreur type} sous l'hypothèse d'indépendance (Lo, 2002, traite le cas général). Avec des rendements non normaux, $B$ change~;
avec des rendements autocorrélés, l'annualisation par $\sqrt{252}$ et la variance changent toutes deux, ce qui fait l'objet de la section suivante~:
la version HAC de la même \omterm{met:qm:estimation:delta}{méthode delta} donne 0.87.

\section{Corrélation sérielle~: variances HAC}

La moyenne d'une série stationnaire, dont les autocovariances sont $\gamma(k)$, a pour variance
\[
\Var(\bar x) = \frac1n\sum_{|k|<n}\Bigl(1 - \frac{|k|}n\Bigr)\gamma(k) \approx \frac1n\sum_k\gamma(k) = \frac{\text{variance de long terme}}n.
\]
Le P\&L du chercheur est la moyenne de cinq positions quotidiennes de cinq jours se chevauchant, si bien que ses autocorrélations valent $1 - k/5$ pour $k
\le 4$ (\cref{fig:qm:estimation:acf}), et la variance de long terme vaut $1 + 2(0.8 + 0.6 + 0.4 + 0.2) = 5$ fois la variance~: l'\omterm{def:qm:estimation:estimator}{erreur type} naïve est
$\sqrt5 = 2.24$ fois trop petite.

\begin{definition}[Estimateur HAC]\label{def:qm:estimation:hac}
Un \omterm{def:qm:estimation:estimator}{estimateur} robuste à l'hétéroscédasticité et à l'autocorrélation (\emph{estimateur HAC}\index{estimateur HAC}) estime une
variance de long terme en pondérant les autocovariances empiriques~; l'\omterm{def:qm:estimation:estimator}{estimateur} de Newey--West $\hat\gamma(0) + 2\sum_{k=1}^L(1 - \frac k{L+1})\hat\gamma(k)$
utilise des poids de Bartlett, qui le rendent non négatif, avec une largeur de bande $L$ croissant lentement avec $n$ (une règle courante est $L =
\lfloor4(n/100)^{2/9}\rfloor$).
\end{definition}

\begin{omfigure}
\begin{tikzpicture}
\begin{groupplot}[group style={group size=2 by 1,horizontal sep=1.6cm},
  width=7.6cm,height=4.8cm,tick label style={font=\small},label style={font=\small},
  grid=major,grid style={gray!25},table/col sep=comma]
\nextgroupplot[xlabel={retard (jours)},ylabel={autocorrélation},xmin=0.3,xmax=10.7,ymin=-0.15,ymax=1,ybar,bar width=6pt,
  legend style={font=\footnotesize,at={(0.97,0.97)},anchor=north east,fill=white,draw=none}]
\addplot[fill=omDef!60,draw=omDef] table[x=lag,y=sample]{figdata/methods/11-estimation/acf.csv};
\addplot[fill=omThm!40,draw=omThm] table[x=lag,y=theory]{figdata/methods/11-estimation/acf.csv};
\legend{échantillon,$1 - k/5$}
\nextgroupplot[xlabel={retards de Newey--West $L$},ylabel={erreur type de la moyenne (pb)},xmin=0,xmax=30,ymin=0,ymax=5,ytick={0,1,2,3},
  legend style={font=\footnotesize,at={(0.97,0.97)},anchor=north east,fill=white,draw=none}]
\addplot[omDef,very thick] table[x=lags,y=nw]{figdata/methods/11-estimation/nwlags.csv};
\addplot[omThm,dashed,thick] table[x=lags,y=iid]{figdata/methods/11-estimation/nwlags.csv};
\addplot[black,dotted,thick] table[x=lags,y=true]{figdata/methods/11-estimation/nwlags.csv};
\legend{Newey--West,formule iid,vraie}
\end{groupplot}
\end{tikzpicture}

\omcaption{Les cinq années de P\&L quotidien du chercheur. À gauche~: autocorrélations empiriques et $1 - k/5$ théorique de cinq
positions se chevauchant. À droite~: \omterm{def:qm:estimation:estimator}{erreur type} de la moyenne du P\&L selon le nombre de retards de Newey--West~: 1.19 pb sans
retard (formule iid), 2.32 aux sept retards de la règle, face à une valeur vraie de 2.65. Données~: le tutoriel du chapitre,
avec graine fixée.}\label{fig:qm:estimation:acf}
\end{omfigure}

Avec sept retards, Newey--West donne 2.32 points de base et une statistique $t$ de 1.73~: le chiffre du collègue. Il reste
inférieur à la vraie valeur de 2.65, car les autocovariances estimées sont biaisées vers zéro dans un échantillon de 1\,260 jours. Sur
2\,000 historiques simulés de la stratégie, les intervalles nominaux à 95~\% construits avec l'\omterm{def:qm:estimation:estimator}{erreur type} iid contiennent la vraie moyenne
dans 62.5~\% des cas~; avec Newey--West, 92.2~\% avec sept retards et 93.5~\% avec vingt (\cref{fig:qm:estimation:coverage}). Une \omterm{def:qm:estimation:estimator}{erreur
type} HAC est nécessaire, mais non suffisante~; l'alternative fondée sur un modèle, lorsque le chevauchement est connu, utilise le facteur 5
directement.

\begin{omfigure}
\begin{tikzpicture}
\begin{axis}[width=9cm,height=4.2cm,xbar,bar width=10pt,
  xlabel={part des intervalles à 95~\% contenant la vraie moyenne (\%)},
  ytick={0,1,2},yticklabels={iid,{NW, 7 retards},{NW, 20 retards}},
  tick label style={font=\small},label style={font=\small},
  xmin=0,xmax=100,ymin=-0.6,ymax=2.6,grid=major,grid style={gray!25},
  table/col sep=comma]
\addplot[fill=omDef!60,draw=omDef,nodes near coords,nodes near coords style={font=\footnotesize,anchor=east,xshift=-1pt}] table[x=coverage,y=k]{figdata/methods/11-estimation/coverage.csv};
\addplot[black,dashed,thick,sharp plot] coordinates {(95,-0.6) (95,2.6)};
\end{axis}
\end{tikzpicture}

\omcaption{Taux de couverture des intervalles de confiance nominaux à 95~\% pour la moyenne du P\&L quotidien de la stratégie se chevauchant, sur 2\,000
historiques simulés de cinq années~: l'\omterm{def:qm:estimation:estimator}{erreur type} iid couvre 62.5~\% du temps, Newey--West 92--94~\%. Données~: le tutoriel du chapitre,
avec graine fixée.}\label{fig:qm:estimation:coverage}
\end{omfigure}

Le même calcul renseigne aussi le chercheur sur le temps à attendre. Un \omterm{def:qm:estimation:sharpe}{ratio de Sharpe} annuel $\mathrm{SR}$ estimé à partir de rendements quotidiens indépendants
atteint une $t$-statistique de 2 après environ $4/\mathrm{SR}^2$ années~: une année pour un \omterm{def:qm:estimation:sharpe}{ratio de Sharpe} de 2, quatre pour 1, seize pour
0.5. Si la variance de long terme vaut $m$ fois la variance, comme pour les positions se chevauchant ($m = 5$), l'annualisation naïve
surestime le \omterm{def:qm:estimation:sharpe}{ratio de Sharpe} d'un facteur $\sqrt m$ et l'attente est $m$ fois plus longue (\cref{fig:qm:estimation:years})~: le 1.51 naïf du chercheur
exige 1.8 années, la lecture honnête 8.8. Cinq années d'historique ne permettent pas de distinguer un \omterm{def:qm:estimation:sharpe}{ratio de Sharpe} honnête de
0.68 de zéro.

\begin{omfigure}
\begin{tikzpicture}
\begin{semilogyaxis}[width=10cm,height=5cm,xlabel={ratio de Sharpe annualisé naïvement},ylabel={années pour atteindre $t = 2$},
  tick label style={font=\small},label style={font=\small},xmin=0.5,xmax=3,ymin=0.3,ymax=100,
  ytick={0.5,1,2,5,10,20,50,100},log ticks with fixed point,grid=major,grid style={gray!25},
  legend columns=3,legend style={font=\footnotesize,at={(0.5,-0.32)},anchor=north,draw=none,fill=none,
  /tikz/every even column/.append style={column sep=8pt}},table/col sep=comma]
\addplot[omDef,very thick] table[x=sr,y=iid]{figdata/methods/11-estimation/years.csv};
\addplot[omThm,very thick,dashed] table[x=sr,y=overlap]{figdata/methods/11-estimation/years.csv};
\addplot[black,dotted,thick,sharp plot] coordinates {(0.5,5) (3,5)};
\addplot[only marks,mark=*,mark size=1.8pt,black] coordinates {(1.51,1.762) (1.51,8.811)};
\legend{rendements indépendants,{chevauchement de cinq jours, $m = 5$},cinq années}
\end{semilogyaxis}
\end{tikzpicture}

\omcaption{Années de données quotidiennes nécessaires avant que le \omterm{def:qm:estimation:sharpe}{ratio de Sharpe} d'une stratégie atteigne une $t$-statistique de 2, en fonction du \omterm{def:qm:estimation:sharpe}{ratio de Sharpe}
annualisé par $\sqrt{252}$~: avec des rendements indépendants et avec les cinq positions se chevauchant du chercheur (variance de long terme
égale à cinq fois la variance). Points~: le 1.51 du chercheur. Données~: le module tutoriel du chapitre.}\label{fig:qm:estimation:years}
\end{omfigure}

\section{Tutoriel~: la $t$-statistique qui a été divisée par deux}\label{tut:qm:estimation:hac}

\begin{tutorial}
\textbf{Objectif.} Reconstruire les \omterm{def:qm:estimation:estimator}{erreurs types} du chercheur et du collègue, mesurer la fréquence à laquelle chacune est juste, et voir le
sandwich corriger une vraisemblance mal spécifiée. \textbf{État final~:}
\cref{fig:qm:estimation:sandwich,fig:qm:estimation:acf,fig:qm:estimation:coverage} et les $t$-statistiques 3.38 et 1.73.
\begin{tutsteps}
\item \textbf{La variance de long terme} avec les poids de Bartlett, et l'\omterm{def:qm:estimation:estimator}{erreur type} d'une moyenne.
\omcode{code/firm/estim/firm_estim.py}{107}{124}{Variance de long terme de Newey--West et erreur type d'une moyenne.}\label{lst:qm:estimation:nw}
\item \textbf{Le \omterm{def:qm:estimation:sharpe}{ratio de Sharpe}} avec ses \omterm{def:qm:estimation:estimator}{erreurs types} iid et HAC, par la \omterm{met:qm:estimation:delta}{méthode delta}.
\omcode{code/firm/estim/firm_estim.py}{127}{145}{Le ratio de Sharpe et ses deux erreurs types.}\label{lst:qm:estimation:sharpe}
\item \textbf{Exécuter} \texttt{problem()}, \texttt{coverage()}, \texttt{misspecified\_fit()} et \texttt{fig\_estimation.py}.
\end{tutsteps}
\textbf{Pour aller plus loin.} Introduire le regroupement de volatilité du P\&L quotidien et comparer les \omterm{def:qm:estimation:estimator}{erreurs types} iid, White et Newey--West~;
estimer la moyenne à partir de rendements de cinq jours non chevauchants et comparer l'efficacité.
\end{tutorial}

\section{Construction~: la boîte à outils d'estimation}\label{bld:qm:estimation:estim}

\begin{build}
\bfield{Rôle} Chaque estimation que la firme miniature publie (une moyenne de rendement, un \omterm{def:qm:estimation:sharpe}{ratio de Sharpe}, un paramètre ajusté) porte une
\omterm{def:qm:estimation:estimator}{erreur type} calculée ici, robuste par défaut.
\bfield{Interface} \texttt{minimize(f, x0)}~; \texttt{gradient}, \texttt{hessian}~; \texttt{mle(loglik\_obs, x0)} renvoyant
l'estimation et ses covariances par le hessien, par le produit externe des scores et par le sandwich~; \texttt{long\_run\_variance(x, lags)}~; \texttt{nw\_lags(n)}~;
\texttt{mean\_se(x, kind)}~; \texttt{sharpe(x, periods)}~; \texttt{delta\_method(g, theta, cov)}.
\bfield{Règles} Des log-vraisemblances par observation (pour pouvoir former les scores)~; le sandwich est rapporté chaque fois qu'une vraisemblance est
maximisée~; largeur de bande HAC selon la règle, sauf indication contraire~; aucune \omterm{def:qm:estimation:estimator}{erreur type} sans que sa méthode soit nommée dans la sortie.
\bfield{Tests d'acceptation} \texttt{code/firm/estim/tests/}~: l'\omterm{def:qm:estimation:mle}{estimateur du maximum de vraisemblance} des échantillons normaux et exponentiels, et leurs \omterm{def:qm:estimation:estimator}{erreurs types}~;
le sandwich égale la variance estimée par le hessien pour un modèle correct~; Newey--West retrouve la variance de long terme d'une moyenne mobile~;
l'\omterm{def:qm:estimation:estimator}{erreur type} de Sharpe concorde avec la simulation.
\bfield{Pour aller plus loin} Le choix automatique de la largeur de bande d'Andrews~; la GMM avec une matrice de pondération efficace et le test de suridentification.
\end{build}

\begin{omsources}
\item P.~J.~Huber, ``The behavior of maximum likelihood estimates under nonstandard conditions'', \emph{Fifth Berkeley
Symposium}, 1967.
\item H.~White, ``Maximum likelihood estimation of misspecified models'', \emph{Econometrica} 50, 1982.
\item W.~K.~Newey et K.~D.~West, ``A simple, positive semi-definite, heteroskedasticity and autocorrelation consistent covariance
matrix'', \emph{Econometrica} 55, 1987.
\item L.~P.~Hansen, ``Large sample properties of generalized method of moments estimators'', \emph{Econometrica} 50, 1982.
\item A.~W.~Lo, ``The statistics of Sharpe ratios'', \emph{Financial Analysts Journal} 58, 2002.
\item G.~Casella et R.~L.~Berger, \emph{Statistical Inference}, Duxbury, 2\up{e} éd., 2002.
\end{omsources}

\section{Exercices}

\begin{exercise}[$\star$]\label{exo:qm:estimation:1}
Cent temps d'attente entre transactions totalisent 50 secondes. Estimer le taux d'arrivée et son \omterm{def:qm:estimation:estimator}{erreur type}.
\end{exercise}

\begin{exercise}[$\star$]\label{exo:qm:estimation:2}
Le \omterm{def:qm:estimation:sharpe}{ratio de Sharpe} annualisé d'une stratégie est estimé à 1.0 à partir de quatre années de rendements quotidiens indépendants. Quelle est son \omterm{def:qm:estimation:estimator}{erreur type}~?
\end{exercise}

\begin{exercise}[$\star$]\label{exo:qm:estimation:3}
La variance échantillonnale de diviseur $n$ est-elle sans biais~? Est-elle convergente~?
\end{exercise}

\begin{exercise}[$\star\star$]\label{exo:qm:estimation:4}
Le P\&L quotidien suit un AR(1) de coefficient 0.2. De quel facteur l'\omterm{def:qm:estimation:estimator}{erreur type} naïve de la moyenne sous-estime-t-elle la
réalité~?
\end{exercise}

\begin{exercise}[$\star\star$]\label{exo:qm:estimation:5}
Un signal a eu raison dans 110 cas sur 200 transactions. En utilisant l'\omterm{def:qm:estimation:mle}{information de Fisher} d'une variable de Bernoulli, donner l'\omterm{def:qm:estimation:estimator}{erreur type} du taux de
réussite, et la $t$-statistique par rapport à 50~\%.
\end{exercise}

\begin{exercise}[$\star\star$]\label{exo:qm:estimation:6}
Pour $n$ observations normales indépendantes, utiliser la \omterm{met:qm:estimation:delta}{méthode delta} pour trouver l'\omterm{def:qm:estimation:estimator}{erreur type} de $\ln\hat\sigma$, et l'évaluer en $n = 2\,000$.
\end{exercise}

\begin{exercise}[$\star\star\star$]\label{exo:qm:estimation:7}
\emph{Programmation.} Sur 2\,000 historiques simulés de la stratégie se chevauchant, mesurer le taux de couverture des intervalles à 95~\% construits sur les
\omterm{def:qm:estimation:estimator}{erreurs types} iid, Newey--West(7) et Newey--West(20).
\end{exercise}

\begin{exercise}[$\star\star\star$]\label{exo:qm:estimation:8}
\emph{Repérer l'erreur.} «~Chaque jour, nous enregistrons le rendement glissant sur 20 jours de notre stratégie. Sur trois ans, le \omterm{def:qm:estimation:sharpe}{ratio de Sharpe} de ces
valeurs quotidiennes, annualisé par $\sqrt{252}$, vaut 1.4, avec une \omterm{def:qm:estimation:estimator}{erreur type} de $1/\sqrt3 = 0.58$~: significatif au seuil de 5~\%.~»
\end{exercise}

\section{Problème~: la $t$-statistique qui a été divisée par deux}

\begin{problem}[{Problème du week-end --- positions qui se chevauchent et erreur type d'une moyenne}]\label{pb:qm:estimation:1}
Une stratégie ouvre une position chaque jour et la conserve cinq jours, de sorte que le P\&L quotidien de chaque jour est la moyenne de cinq positions qui se chevauchent.
Sur cinq ans (1\,260 jours), son P\&L quotidien, exprimé en points de base du capital, est l'historique simulé du chapitre.

\medskip\noindent\textbf{Partie I --- La vision naïve.}
\begin{enumerate}
\item Quelles sont la moyenne et l'écart type empiriques~?
\item Quelles sont l'\omterm{def:qm:estimation:estimator}{erreur type} et la $t$-statistique sous hypothèse iid~?
\item Quels sont le \omterm{def:qm:estimation:sharpe}{ratio de Sharpe} annualisé et son \omterm{def:qm:estimation:estimator}{erreur type} iid~?
\item Quelles sont les quatre premières autocorrélations empiriques, comparées à la théorie~?
\item Pourquoi les jours consécutifs sont-ils corrélés~?
\end{enumerate}

\medskip\noindent\textbf{Partie II --- \omterm{def:qm:estimation:estimator}{Erreurs types} robustes.}
\begin{enumerate}[resume]
\item Quel est le facteur de variance de long terme du P\&L, et quelle est la vraie \omterm{def:qm:estimation:estimator}{erreur type} de la moyenne~?
\item Combien de retards la règle donne-t-elle, et que deviennent Newey--West et la $t$-statistique~?
\item Et avec vingt retards~?
\item Quelle est l'\omterm{def:qm:estimation:estimator}{erreur type} HAC du \omterm{def:qm:estimation:sharpe}{ratio de Sharpe}~?
\item À quelle fréquence les intervalles iid à 95~\% contiennent-ils la vraie moyenne, sur 2\,000 historiques simulés~?
\end{enumerate}

\medskip\noindent\textbf{Partie III --- Ce que les données permettent de dire.}
\begin{enumerate}[resume]
\item Connaissant le chevauchement, quelle \omterm{def:qm:estimation:estimator}{erreur type} fondée sur un modèle utiliseriez-vous~?
\item Quelle $t$-statistique le chercheur doit-il attendre en moyenne, avec une vraie moyenne de 4~pb~?
\item Combien d'années faudrait-il en moyenne pour obtenir une $t$-statistique de 3~?
\item Pourquoi Newey--West reste-t-il légèrement trop petit~?
\item Que changerait l'évaluation du P\&L sur cinq jours non chevauchants~?
\end{enumerate}

\medskip\noindent\textbf{Partie IV --- Jugement.}
\begin{enumerate}[resume]
\item Quelle autre dépendance du P\&L quotidien la formule iid manquerait-elle~?
\item Le chercheur a testé vingt variantes avant celle-ci. Que faut-il changer d'autre dans l'évaluation~?
\item Que doit indiquer le rapport sur la stratégie~?
\item Énoncer le \emph{résultat nommé}~: les $t$-statistiques naïve et robuste.
\item En une phrase~: quand l'\omterm{def:qm:estimation:estimator}{erreur type} iid d'une moyenne est-elle fausse~?
\end{enumerate}
\end{problem}

\section{Questions d'entretien}

\begin{interviewq}[$\star$ \iqroles{chercheur, ingénieur ML}]\label{iq:qm:estimation:1}
Qu'est-ce que le maximum de vraisemblance, et pourquoi est-ce l'\omterm{def:qm:estimation:estimator}{estimateur} par défaut~?
\end{interviewq}

\begin{interviewq}[$\star\star$ \iqroles{chercheur, trader}]\label{iq:qm:estimation:2}
Une stratégie affiche un \omterm{def:qm:estimation:sharpe}{ratio de Sharpe} de 1.0 sur un an. Quelle confiance accordez-vous au fait qu'il soit positif~?
\end{interviewq}

\begin{interviewq}[$\star\star$ \iqroles{chercheur}]\label{iq:qm:estimation:3}
Qu'est-ce qu'une \omterm{def:qm:estimation:sandwich}{variance sandwich}, et quand en avez-vous besoin~?
\end{interviewq}

\begin{interviewq}[$\star\star$ \iqroles{chercheur}]\label{iq:qm:estimation:4}
Vous calculez chaque jour des rendements mensuels à partir de fenêtres qui se chevauchent, puis vous les régressez sur un signal. Qu'est-ce qui ne va pas avec les \omterm{def:qm:estimation:estimator}{erreurs types} habituelles, et comment les corrigez-vous~?
standard errors, and how do you fix them~?
\end{interviewq}

\begin{interviewq}[$\star\star$ \iqroles{chercheur, ingénieur ML}]\label{iq:qm:estimation:5}
Qu'est-ce que l'\omterm{def:qm:estimation:mle}{information de Fisher}, et que dit la borne de Cramér--Rao~?
\end{interviewq}

\begin{interviewq}[$\star\star\star$ \iqroles{chercheur}]\label{iq:qm:estimation:6}
Si votre vraisemblance est fausse, qu'estime son maximiseur~?
\end{interviewq}
